\documentclass[a4paper,12pt]{ctexart}
\usepackage{../teaching-style}

\begin{document}
\pagestyle{empty}

\maintitle{计算进阶：式的运算}
\begin{center}
  \large 适用于5\~{}初一 \quad 带讲+答案
\end{center}
\vspace{0.3cm}

\chaptertitle{第一讲\quad 整式加减乘除}

\sectiontitle{知识梳理}
\knowbox{
\textbf{1. 整式的概念}\\
单项式：由数与字母的积组成的代数式。单项式的系数、次数。\\
多项式：几个单项式的和。多项式的项、次数。\\
\textbf{2. 整式加减}\\
去括号、合并同类项。\\
同类项：字母相同且相同字母的指数也相同的项。\\
\textbf{3. 整式乘法}\\
单项式$\times$单项式：系数相乘，同底数幂相乘。\\
单项式$\times$多项式：分配律。\\
多项式$\times$多项式：逐项相乘。\\
\textbf{4. 整式除法}\\
单项式$\div$单项式：系数相除，同底数幂相除。\\
多项式$\div$单项式：逐项相除。}

\sectiontitle{典型例题}

\begin{example}
计算：$(2x^2-3xy+4y^2)-2(x^2-xy+2y^2)$
\end{example}
\begin{solution}
原式 $=2x^2-3xy+4y^2-2x^2+2xy-4y^2=-xy$
\end{solution}

\begin{example}
计算：$(2a-3b)(3a+2b)$
\end{example}
\begin{solution}
原式 $=2a\cdot3a+2a\cdot2b-3b\cdot3a-3b\cdot2b=6a^2+4ab-9ab-6b^2=6a^2-5ab-6b^2$
\end{solution}

\begin{example}
计算：$(6x^3y^2-9x^2y^3+12xy^4)\div(3xy^2)$
\end{example}
\begin{solution}
原式 $=2x^2-3xy+4y^2$
\end{solution}

\sectiontitle{分层练习}

\begin{exercise}
\item 计算：
  \begin{subexercise}
    \item $(3x^2-2x+1)+(x^2+3x-5)$
    \item $(5a^2-3ab+2b^2)-(2a^2+ab-b^2)$
    \item $2(3x-2y)-3(x-3y)$
    \item $3(2x^2-xy)-2(3x^2+xy-1)$
  \end{subexercise}

\item 计算：
  \begin{subexercise}
    \item $(-2a^2b^3)\times(-3a^3b^2)$
    \item $4x^2y\times(-3xy^2)^2$
    \item $(2x+1)(x-3)$
    \item $(3a-2b)(2a+3b)$
    \item $(x+2y)(x-3y)$
    \item $(2m-3n)(3m+2n)$
  \end{subexercise}

\item 计算：
  \begin{subexercise}
    \item $(8a^3b^2-12a^2b^3)\div(4a^2b)$
    \item $(15x^3y^2-10x^2y^3+5x^2y^2)\div(5x^2y)$
    \item $(12m^3n^2-18m^2n^3+6m^2n^2)\div(6m^2n)$
    \item $(9x^4y^3-12x^3y^4+6x^3y^3)\div(3x^3y^2)$
  \end{subexercise}

\item 先化简，再求值：
  \begin{subexercise}
    \item $(2x-3y)+(5x+4y)$，其中$x=-1$，$y=2$
    \item $(3a^2-ab+2b^2)-(a^2+ab-2b^2)$，其中$a=2$，$b=-1$
  \end{subexercise}
\end{exercise}

\newpage

\chaptertitle{第二讲\quad 乘法公式}

\sectiontitle{知识梳理}
\knowbox{
\textbf{1. 平方差公式}\\
$(a+b)(a-b)=a^2-b^2$\\
\textbf{2. 完全平方公式}\\
$(a+b)^2=a^2+2ab+b^2$\\
$(a-b)^2=a^2-2ab+b^2$\\
\textbf{3. 立方和与立方差}\\
$(a+b)(a^2-ab+b^2)=a^3+b^3$\\
$(a-b)(a^2+ab+b^2)=a^3-b^3$\\
\textbf{4. 三个数的完全平方}\\
$(a+b+c)^2=a^2+b^2+c^2+2ab+2bc+2ca$}

\sectiontitle{典型例题}

\begin{example}
计算：$(3x+2y)(3x-2y)$
\end{example}
\begin{solution}
原式 $=(3x)^2-(2y)^2=9x^2-4y^2$
\end{solution}

\begin{example}
计算：$(2x-3y)^2$
\end{example}
\begin{solution}
原式 $=(2x)^2-2\cdot(2x)\cdot(3y)+(3y)^2=4x^2-12xy+9y^2$
\end{solution}

\begin{example}
计算：$(x+2)(x-2)(x^2+4)$
\end{example}
\begin{solution}
原式 $=(x^2-4)(x^2+4)=x^4-16$
\end{solution}

\sectiontitle{分层练习}

\begin{exercise}
\item 利用平方差公式计算：
  \begin{subexercise}
    \item $(x+5)(x-5)$
    \item $(2a+3b)(2a-3b)$
    \item $(-3x+2y)(-3x-2y)$
    \item $(m+2n)(m-2n)$
    \item $(5x+1)(5x-1)-25x^2$
    \item $(3a-2b)(3a+2b)-(2a+3b)(2a-3b)$
  \end{subexercise}

\item 利用完全平方公式计算：
  \begin{subexercise}
    \item $(x+4)^2$
    \item $(2a-3b)^2$
    \item $(-3x+2y)^2$
    \item $(m-5n)^2$
    \item $(3x+2y)^2-(3x-2y)^2$
    \item $(a+b)^2-(a-b)^2$
  \end{subexercise}

\item 计算：
  \begin{subexercise}
    \item $(x+1)(x-1)(x^2+1)$
    \item $(x+2)(x^2-2x+4)$
    \item $(2a-1)(4a^2+2a+1)$
    \item $(x+2)^2-(x-2)^2$
    \item $(a+b+c)^2$
    \item $(2x-y+3z)^2$
  \end{subexercise}

\item 简便运算：
  \begin{subexercise}
    \item $102\times98$
    \item $199^2$
    \item $1002\times998$
    \item $201^2-199^2$
  \end{subexercise}
\end{exercise}

\newpage

\chaptertitle{第三讲\quad 因式分解（一）}

\sectiontitle{知识梳理}
\knowbox{
\textbf{1. 因式分解的概念}\\
把一个多项式化成几个整式的积的形式，叫做因式分解。\\
\textbf{2. 提公因式法}\\
$ma+mb+mc=m(a+b+c)$\\
\textbf{3. 公式法}\\
平方差公式：$a^2-b^2=(a+b)(a-b)$\\
完全平方公式：$a^2\pm2ab+b^2=(a\pm b)^2$\\
立方和/差：$a^3\pm b^3=(a\pm b)(a^2\mp ab+b^2)$\\
\textbf{4. 十字相乘法}\\
$x^2+(p+q)x+pq=(x+p)(x+q)$\\
$acx^2+(ad+bc)x+bd=(ax+b)(cx+d)$}

\sectiontitle{典型例题}

\begin{example}
分解因式：$12a^3b^2-18a^2b^3+24a^2b^2$
\end{example}
\begin{solution}
原式 $=6a^2b^2(2a-3b+4)$
\end{solution}

\begin{example}
分解因式：$4x^2-9y^2$
\end{example}
\begin{solution}
原式 $=(2x)^2-(3y)^2=(2x+3y)(2x-3y)$
\end{solution}

\begin{example}
分解因式：$x^2-5x+6$
\end{example}
\begin{solution}
原式 $=(x-2)(x-3)$
\end{solution}

\sectiontitle{分层练习}

\begin{exercise}
\item 提公因式：
  \begin{subexercise}
    \item $6a^2b-9ab^2$
    \item $15x^3y^2-10x^2y^3+5x^2y$
    \item $3m^3n-12mn^3+6m^2n^2$
    \item $-8a^3b^2+12a^2b^3-4a^2b^2$
  \end{subexercise}

\item 用平方差公式分解：
  \begin{subexercise}
    \item $x^2-16$
    \item $4a^2-9b^2$
    \item $25x^2-1$
    \item $9m^2-16n^2$
    \item $x^4-y^4$
    \item $a^2b^2-1$
  \end{subexercise}

\item 用完全平方公式分解：
  \begin{subexercise}
    \item $x^2+6x+9$
    \item $4a^2-20a+25$
    \item $9x^2+12xy+4y^2$
    \item $m^2-8mn+16n^2$
    \item $x^2y^2-2xy+1$
    \item $a^4-2a^2b^2+b^4$
  \end{subexercise}

\item 十字相乘法：
  \begin{subexercise}
    \item $x^2+7x+12$
    \item $x^2-5x-14$
    \item $x^2+2x-15$
    \item $x^2-8x+15$
    \item $2x^2+7x+3$
    \item $3x^2-10x+3$
  \end{subexercise}
\end{exercise}

\newpage

\chaptertitle{第四讲\quad 分式四则运算}

\sectiontitle{知识梳理}
\knowbox{
\textbf{1. 分式的概念}\\
$\dfrac{A}{B}$（$B$中含有字母）叫做分式。分母不为零。\\
\textbf{2. 分式基本性质}\\
$\dfrac{A}{B}=\dfrac{A\times C}{B\times C}$，$\dfrac{A}{B}=\dfrac{A\div C}{B\div C}$（$C\neq0$）\\
\textbf{3. 分式的运算}\\
加减：$\dfrac{a}{b}\pm\dfrac{c}{d}=\dfrac{ad\pm bc}{bd}$\\
乘法：$\dfrac{a}{b}\times\dfrac{c}{d}=\dfrac{ac}{bd}$\\
除法：$\dfrac{a}{b}\div\dfrac{c}{d}=\dfrac{a}{b}\times\dfrac{d}{c}=\dfrac{ad}{bc}$\\
乘方：$\left(\dfrac{a}{b}\right)^n=\dfrac{a^n}{b^n}$}

\sectiontitle{典型例题}

\begin{example}
化简：$\dfrac{x^2-4}{x^2+4x+4}$
\end{example}
\begin{solution}
原式 $=\dfrac{(x+2)(x-2)}{(x+2)^2}=\dfrac{x-2}{x+2}$
\end{solution}

\begin{example}
计算：$\dfrac{3}{a}-\dfrac{2}{a^2}+\dfrac{1}{3a}$
\end{example}
\begin{solution}
原式 $=\dfrac{9a}{3a^2}-\dfrac{6}{3a^2}+\dfrac{a}{3a^2}=\dfrac{9a-6+a}{3a^2}=\dfrac{10a-6}{3a^2}$
\end{solution}

\sectiontitle{分层练习}

\begin{exercise}
\item 化简：
  \begin{subexercise}
    \item $\dfrac{x^2-1}{x^2+2x+1}$
    \item $\dfrac{a^2-4}{a^2-4a+4}$
    \item $\dfrac{6x^2y^3}{9xy^2}$
    \item $\dfrac{m^2-9}{m^2-6m+9}$
  \end{subexercise}

\item 计算：
  \begin{subexercise}
    \item $\dfrac{2}{x}+\dfrac{3}{x^2}$
    \item $\dfrac{5}{ab}-\dfrac{3}{bc}$
    \item $\dfrac{1}{a-1}+\dfrac{1}{a+1}$
    \item $\dfrac{2}{x-3}-\dfrac{1}{x+3}$
  \end{subexercise}

\item 计算：
  \begin{subexercise}
    \item $\dfrac{a}{b^2}\times\dfrac{b}{a^2}$
    \item $\dfrac{x^2-4}{x+3}\div\dfrac{x-2}{x+3}$
    \item $\left(\dfrac{a}{b}\right)^3\times\left(\dfrac{b^2}{a}\right)^2$
    \item $\dfrac{x^2-1}{x^2-4}\div\dfrac{x-1}{x+2}$
  \end{subexercise}

\item 先化简，再求值：$\dfrac{x^2-1}{x^2+2x+1}\div\dfrac{x-1}{x+1}$，其中$x=2$
\end{exercise}

\newpage

\chaptertitle{第五讲\quad 二次根式化简与运算}

\sectiontitle{知识梳理}
\knowbox{
\textbf{1. 二次根式}\\
形如$\sqrt{a}(a\ge0)$的式子叫做二次根式。\\
\textbf{2. 性质}\\
$\sqrt{a^2}=|a|=\begin{cases}a&(a\ge0)\\-a&(a<0)\end{cases}$\\
$(\sqrt{a})^2=a\;(a\ge0)$\\
$\sqrt{a}\cdot\sqrt{b}=\sqrt{ab}\;(a,b\ge0)$\\
$\dfrac{\sqrt{a}}{\sqrt{b}}=\sqrt{\dfrac{a}{b}}\;(a\ge0,b>0)$\\
\textbf{3. 最简二次根式}\\
被开方数不含分母，不含能开得尽方的因数或因式。\\
\textbf{4. 分母有理化}\\
$\dfrac{1}{\sqrt{a}}=\dfrac{\sqrt{a}}{a}$，$\dfrac{1}{\sqrt{a}+\sqrt{b}}=\dfrac{\sqrt{a}-\sqrt{b}}{a-b}$}

\sectiontitle{典型例题}

\begin{example}
化简：$\sqrt{48}-\sqrt{75}+\sqrt{27}$
\end{example}
\begin{solution}
原式 $=4\sqrt3-5\sqrt3+3\sqrt3=2\sqrt3$
\end{solution}

\begin{example}
计算：$\dfrac{\sqrt{15}+\sqrt{5}}{\sqrt5}$
\end{example}
\begin{solution}
原式 $=\dfrac{\sqrt{15}}{\sqrt5}+\dfrac{\sqrt5}{\sqrt5}=\sqrt3+1$
\end{solution}

\sectiontitle{分层练习}

\begin{exercise}
\item 化简：
  \begin{subexercise}
    \item $\sqrt{32}$
    \item $\sqrt{72}$
    \item $\sqrt{98}$
    \item $\sqrt{108}$
    \item $\sqrt{\dfrac{1}{2}}$
    \item $\sqrt{\dfrac{3}{4}}$
  \end{subexercise}

\item 计算：
  \begin{subexercise}
    \item $\sqrt{8}+\sqrt{18}-\sqrt{32}$
    \item $\sqrt{27}+2\sqrt{\dfrac{1}{3}}-\sqrt{48}$
    \item $\sqrt{50}-3\sqrt{2}+\sqrt{8}$
    \item $\sqrt{75}+\sqrt{\dfrac{1}{3}}-\sqrt{12}$
  \end{subexercise}

\item 分母有理化：
  \begin{subexercise}
    \item $\dfrac{1}{\sqrt3}$
    \item $\dfrac{2}{\sqrt5}$
    \item $\dfrac{3}{\sqrt{12}}$
    \item $\dfrac{1}{\sqrt2-1}$
    \item $\dfrac{3}{\sqrt5+\sqrt2}$
    \item $\dfrac{2}{\sqrt3-\sqrt2}$
  \end{subexercise}
\end{exercise}

\newpage

\chaptertitle{综合闯关}

\begin{exercise}
\item 计算：$(3x^2-2x+5)-2(x^2-3x-2)$
\item 计算：$(2a-3b)^2-(a+2b)(a-2b)$
\item 简便计算：$103\times97$
\item 分解因式：$x^3-4x$
\item 分解因式：$x^2-6x+9-4y^2$
\item 分解因式：$2x^2-7x+3$
\item 计算：$\dfrac{x+1}{x-1}-\dfrac{x^2+1}{x^2-1}$
\item 计算：$\left(\dfrac{x}{y}-\dfrac{y}{x}\right)\div\dfrac{x-y}{x}$
\item 化简：$\sqrt{48}-\sqrt{27}+\sqrt{\dfrac{4}{3}}$
\item 分母有理化：$\dfrac{2}{\sqrt5+\sqrt3}+\dfrac{1}{\sqrt3-1}$
\end{exercise}

\newpage

\chaptertitle{参考答案}

\textbf{第一讲}
\noindent\\
1.(a) $4x^2+x-4$ \quad (b) $3a^2-4ab+3b^2$ \quad (c) $3x+5y$ \quad (d) $-5xy+2$\\
2.(a) $6a^5b^5$ \quad (b) $36x^4y^5$ \quad (c) $2x^2-5x-3$ \quad (d) $6a^2+5ab-6b^2$\\
(e) $x^2-xy-6y^2$ \quad (f) $6m^2-5mn-6n^2$\\
3.(a) $2ab-3b^2$ \quad (b) $3xy-2y^2+y$ \quad (c) $2mn-3n^2+n$ \quad (d) $3xy-4y^2+2y$\\
4.(a) $2$ \quad (b) $14$\\

\textbf{第二讲}
\noindent\\
1.(a) $x^2-25$ \quad (b) $4a^2-9b^2$ \quad (c) $9x^2-4y^2$ \quad (d) $m^2-4n^2$\\
(e) $-1$ \quad (f) $5a^2-5b^2$\\
2.(a) $x^2+8x+16$ \quad (b) $4a^2-12ab+9b^2$ \quad (c) $9x^2-12xy+4y^2$\\
(d) $m^2-10mn+25n^2$ \quad (e) $24xy$ \quad (f) $4ab$\\
3.(a) $x^4-1$ \quad (b) $x^3+8$ \quad (c) $8a^3-1$ \quad (d) $8x$ \quad (e) $a^2+b^2+c^2+2ab+2bc+2ca$\\
(f) $4x^2+y^2+9z^2-4xy+12xz-6yz$\\
4.(a) $9996$ \quad (b) $39601$ \quad (c) $999996$ \quad (d) $800$\\

\textbf{第三讲}
\noindent\\
1.(a) $3ab(2a-3b)$ \quad (b) $5x^2y(3xy-2y^2+1)$ \quad (c) $3mn(m^2-4n^2+2mn)$ \quad (d) $-4a^2b^2(2a-3b+1)$\\
2.(a) $(x+4)(x-4)$ \quad (b) $(2a+3b)(2a-3b)$ \quad (c) $(5x+1)(5x-1)$ \quad (d) $(3m+4n)(3m-4n)$\\
(e) $(x^2+y^2)(x+y)(x-y)$ \quad (f) $(ab+1)(ab-1)$\\
3.(a) $(x+3)^2$ \quad (b) $(2a-5)^2$ \quad (c) $(3x+2y)^2$ \quad (d) $(m-4n)^2$\\
(e) $(xy-1)^2$ \quad (f) $(a^2-b^2)^2$\\
4.(a) $(x+3)(x+4)$ \quad (b) $(x-7)(x+2)$ \quad (c) $(x+5)(x-3)$ \quad (d) $(x-3)(x-5)$\\
(e) $(2x+1)(x+3)$ \quad (f) $(3x-1)(x-3)$\\

\textbf{第四讲}
\noindent\\
1.(a) $\dfrac{x-1}{x+1}$ \quad (b) $\dfrac{a+2}{a-2}$ \quad (c) $\dfrac{2xy}{3}$ \quad (d) $\dfrac{m+3}{m-3}$\\
2.(a) $\dfrac{2x+3}{x^2}$ \quad (b) $\dfrac{5c-3a}{abc}$ \quad (c) $\dfrac{2a}{a^2-1}$ \quad (d) $\dfrac{x+9}{x^2-9}$\\
3.(a) $\dfrac{1}{ab}$ \quad (b) $x+2$ \quad (c) $a$ \quad (d) $\dfrac{x+1}{x-2}$\\
4. $1$\\

\textbf{第五讲}
\noindent\\
1.(a) $4\sqrt2$ \quad (b) $6\sqrt2$ \quad (c) $7\sqrt2$ \quad (d) $6\sqrt3$ \quad (e) $\dfrac{\sqrt2}{2}$ \quad (f) $\dfrac{\sqrt3}{2}$\\
2.(a) $\sqrt2$ \quad (b) $\dfrac{4\sqrt3}{3}$ \quad (c) $4\sqrt2$ \quad (d) $\dfrac{14\sqrt3}{3}$\\
3.(a) $\dfrac{\sqrt3}{3}$ \quad (b) $\dfrac{2\sqrt5}{5}$ \quad (c) $\dfrac{\sqrt3}{2}$ \quad (d) $\sqrt2+1$ \quad (e) $\sqrt5-\sqrt2$ \quad (f) $2\sqrt3+2\sqrt2$\\

\textbf{综合闯关}
\noindent\\
1. $x^2+4x+9$ \quad 2. $3a^2-12ab+13b^2$ \quad 3. $9991$ \quad 4. $x(x+2)(x-2)$ \quad 5. $(x-3+2y)(x-3-2y)$\\
6. $(2x-1)(x-3)$ \quad 7. $\dfrac{-2}{x+1}$ \quad 8. $\dfrac{x+y}{y}$ \quad 9. $\dfrac{5\sqrt3}{3}$ \quad 10. $\sqrt5$

\end{document}
